\begin{abstract}

We show that quadratic $f(R)$ gravity models with regularized mass profiles exhibit a structurally controlled critical transition in horizon formation. This transition is governed by a dimensionless parameter $\beta$ and separates horizon-bearing from horizonless compact configurations.

The critical behavior originates from the root structure of the horizon equation and is associated with a non-analytic collapse of the outer horizon branch at a critical value $\beta_c$. This transition corresponds to a saddle-node bifurcation, indicating that the existence of horizons is controlled by a structural condition rather than continuous deformation.

Near the critical regime, the effective metric develops an extended near-zero region, leading to a divergence of the echo delay time and the formation of an effectively infinite optical cavity. This behavior is analogous to critical slowing down in dynamical systems and provides a distinct physical signature of the transition.

These results indicate that critical horizon transitions may represent a generic feature of regular compact-object geometries and provide observable signatures in shadow structure, quasi-normal mode spectra, and gravitational-wave echo signals.

\end{abstract}
